\section{Discussion}
\label{sec:discussion}

\subsection{Internal spin structure and momentum-resolved weight}

The projected response in Eq.~\eqref{eq:folded-response} separates two
physical ingredients. The effective chain determines the collective spectrum
at cell momentum $Q$, while the five-site spin distribution determines its
visibility at physical momentum $q$ through $F(q)/5$.
Consequently, a weak response at one replica need not indicate a missing
collective mode. At negative $\alpha$, the near cancellation of the
$q=\pi$ cell amplitude coexists with a strong amplitude at $q=0.6\pi$,
although both momenta fold to $Q=\pi$.
At positive $\alpha$ the same comparison favors $q=\pi$.
The switch of the isolated cell from spin $1/2$ to spin $3/2$ changes this
interference pattern as well as the scale of the inter-cell exchange.

The two-cell calculation tests this statement without a phenomenological
line shape or an intensity fit. At $\alpha=\pm0.5$ and
$\beta/\Delta_{\mathrm{loc}}\le0.01$, the largest relative error of the
41-momentum low-band weight vector is $0.27\%$, and that of the
singlet--triplet gap is $0.56\%$.
These checks validate the absolute projected weights for the specified
finite system in the weak-coupling limit.
For longer chains the same leading factor $F(q)/5$ multiplies the effective
response, but the two-cell line spectrum and its integrated absolute
weights cannot be substituted for those of an extended chain.

This use of the cluster projection complements the trimer and tetramer
studies, which compared numerical spectra with effective descriptions of
propagating excitations~\cite{cheng2022trimer,li2026tetramer}.
Here the quantitative comparison concerns the low-energy intensity associated
with two different half-integer cell spins and their folded momentum
replicas. It does not identify the high-energy bands of
Fig.~\ref{fig:fullq} with unique particles.
Indeed, the two plotted points have
$\beta/\Delta_{\mathrm{loc}}=1.46$ and $3.96$ for positive and negative
$\alpha$, respectively. The two-cell controls at these points already fail
the separation required by the weak-coupling low-band construction.
Their $L=40$ spectra therefore constitute a finite-coupling comparison,
not an extrapolation of the absolute weak-coupling formula.

\subsection{A different local limit near zero intra-cell exchange}

The quartet projection for negative $\alpha$ and the reconstruction toward
$q=\pi$ near $\alpha=0$ are not contradictory: the projection presumes that
the isolated-cell gap is large compared with $\beta$.
At fixed positive $\beta$ this condition fails as $\alpha$ approaches zero,
and a different partition of the chain becomes appropriate.
In the bulk at $\alpha=0$, sites 2--4 of each cell form a spin-$1/2$ trimer,
while the bond joining site 5 to site 1 of the next cell forms a singlet
dimer. Open boundaries additionally leave uncoupled end spins at this
decoupled point.

For $|\alpha|\ll\min(1,\beta)$ the trimer doublets interact through a virtual
triplet of the intervening dimer. Each trimer endpoint projects onto
$2\bm T/3$, and eliminating the dimer triplet yields
\begin{equation}
\Jeff=\frac{2\alpha^2}{9\beta}>0
\label{eq:near-zero}
\end{equation}
in units of $J_1$.
Thus both signs of sufficiently small $\alpha$ lead, at this order, to
antiferromagnetically coupled trimer doublets.
In an open chain, however, each unpaired end spin couples to its adjacent
trimer already at first order, as
$(2\alpha/3)\bm S_{\mathrm{edge}}\!\cdot\!\bm T$.
This sign-dependent $O(\alpha)$ boundary term is larger than the bulk
$O(\alpha^2/\beta)$ scale in the stated limit.
Equation~\eqref{eq:near-zero} is therefore not a formula for the lowest
sector gap of the complete open chain.
Exact diagonalization of the eight-spin trimer--dimer--trimer cluster
supports the coefficient and the $\alpha^2/\beta$ scaling in the
asymptotic limit. The corresponding local-gap estimates at
$(\alpha,\beta)=(-0.10,0.25)$ and $(-0.175,0.25)$ are already
$36.1\%$ and $53.9\%$ below the leading prediction, respectively.
Equation~\eqref{eq:near-zero} therefore explains the limiting mechanism
and its small energy scale, not the finite-coupling location of the
momentum reconstruction.

The static scan compares these competing local descriptions in finite
chains. Its dominant momentum changes from about $0.6\pi$ toward $\pi$,
but the momentum change and the size trend of the peak sharpness do not
select the same coupling intervals.
A change in the strongest physical-momentum replica alone is not a
thermodynamic phase criterion: it can reflect changes of matrix elements
and short-distance correlations as well as changes of long-distance order.
The data establish the finite-size reconstruction but do not determine
whether any part of it becomes singular in the bulk limit.
The small scale in Eq.~\eqref{eq:near-zero} also explains why a low absolute
energy variance need not by itself control the entropy or a sector gap near
this limit; the solver-sensitivity checks must be judged against the relevant
low-energy scale.

\subsection{What a spin probe can distinguish}

There are two separate uses of interference in the present analysis.
Within the cell ground multiplet, $F(q)$ predicts the relative low-energy
weights of folded replicas. Between different isolated-cell multiplets,
reflection parity can forbid a transition exactly.
For the positive-$\alpha$ pentamer, the $q=\pi$ spin operator is even under
reflection, so the lowest odd-parity excited doublet is dark and higher
allowed sectors carry the local inelastic weight.
Coupling the cells does not conserve the reflection of each cell
independently. The isolated-cell dark state is therefore a weak-coupling
selection rule to be followed under mixing, not an exact prohibition of an
entire many-body band at arbitrary $\beta$.

These distinctions suggest direct tests of a weakly coupled pentamer
realization: compare low-energy intensities at momenta sharing the same
$Q$, and compare the lowest local excitation energy with the first allowed
single-spin signal. Quantitative comparison with a material additionally
requires its positions, exchange constants, and scattering operator.
For an SU(2)-invariant singlet the three Cartesian spin responses are equal,
so $S^{zz}$ determines the spin-correlation tensor before the experimental
form factors and polarization factors are applied.
The common finite-time resolution of Fig.~\ref{fig:fullq},
$0.09277J_1$, is broader than the separately calculated finite-chain gaps.
The maps compare resolved intensity distributions; they do not resolve
those gaps from zero or establish the asymptotic widths of spectral features.
